\section{Conclusion}
\label{sec:conclusion}

We presented \dgtdm, a three-branch framework for fine-grained action
recognition. The task is hard because the classes look almost identical in
appearance, and the difference lies in how the body moves. We decomposed this
into two problems and gave each to a separate branch. A frozen CLIP encoder
handles localization, supplying stable spatial priors learned from
large-scale pre-training, and a spatio-temporal guidance branch uses them to
focus a lightweight temporal encoder on the action region. Because a 2D
backbone has no temporal modeling of its own, a dual-granularity temporal
dynamics module adds it on the 2D feature stream: MoDeA tracks short-term
body-part motion with dense deformable sampling, and TSA keeps a long-term
memory over the sequence to characterize the overall temporal response. The
framework thus targets the two difficulties that make fine-grained actions
hard, namely similar backgrounds and subtle body movement, without inflating
the backbone into a 3D network.

Ablations on Diving48 confirm that the two temporal branches and the guidance
all carry part of the accuracy: every variant that removes the deformable
sampling of MoDeA, turns the TSA state into a forward-only scan, or disables
the per-stage prior injection drops below the full model. The frozen CLIP
expert could in principle be replaced by a cheaper spatial encoder trained
end-to-end; we leave that exchange, together with an extension to further
benchmarks, to future work.
